\documentclass[11pt,a4paper]{report}

% Packages for technical specification
\usepackage[margin=0.85in]{geometry}
\usepackage{fancyhdr}
\usepackage{titlesec}
\usepackage{booktabs}
\usepackage{longtable}
\usepackage{array}
\usepackage{amsmath,amssymb}
\usepackage{enumitem}
\usepackage{hyperref}
\usepackage{xcolor}
\usepackage{graphicx}
\usepackage{float}
\usepackage{verbatim}
\usepackage{listings}
\usepackage{tikz}
\usetikzlibrary{shapes.geometric, arrows.meta, positioning, calc, fit, backgrounds, shadows}

% Colors
\definecolor{urteg}{RGB}{0,128,128}
\definecolor{posixblue}{RGB}{0,51,102}
\definecolor{kernelgreen}{RGB}{0,100,0}
\definecolor{sectblue}{RGB}{0,51,102}

% Listings style for code/schemas
\lstdefinestyle{kernelschema}{
    basicstyle=\ttfamily\footnotesize,
    breaklines=true,
    frame=single,
    backgroundcolor=\color{gray!6},
    keywordstyle=\color{posixblue}\bfseries,
    commentstyle=\color{gray},
    numbers=left,
    numberstyle=\tiny\color{gray}
}

% Header/Footer
\pagestyle{fancy}
\fancyhf{}
\fancyhead[L]{\small URTE POSIX Kernel Technical Specification | Confidential}
\fancyhead[R]{\small June 2026 | v1.0 | IEEE Std 1003.1 Compliant}
\fancyfoot[C]{\thepage}
\renewcommand{\headrulewidth}{0.4pt}

% Section formatting
\titleformat{\chapter}{\Large\bfseries\color{sectblue}}{\thechapter.}{0.5em}{}
\titleformat{\section}{\large\bfseries\color{posixblue}}{\thesection}{0.5em}{}
\titleformat{\subsection}{\normalsize\bfseries\color{urteg}}{\thesubsection}{0.5em}{}

% Title
\title{\textbf{URTE POSIX-Compliant Operating System Kernel}\\[0.3cm]
\large Technical Specification\\[0.2cm]
\normalsize IEEE Std 1003.1 (POSIX) \& Single UNIX Specification Compliant\\[0.3cm]
for the\\
\textbf{Universal Regeneration Therapy Ecosystem (URTE)}}
\author{
    \textbf{Prepared by:}\\[0.1cm]
    Grok xAI Research \& UAB Saptera\\[0.1cm]
    Kristupas Gelzinis \& Olek Suchodolski\\[0.3cm]
    \textit{Based on URTE Architecture Document, MBSE Baseline,}\\
    \textit{Arcadia/Capella Extension, and Infrastructure Model (June 2026)}
}
\date{June 2026}

\begin{document}

\maketitle

\begin{abstract}
\noindent This document provides an \textbf{exhaustive technical specification} for the \textbf{URTE Operating System Kernel} — the foundational runtime environment for the Universal Regeneration Therapy Ecosystem.

The kernel is designed to be fully \textbf{POSIX-compliant} (IEEE Std 1003.1-2024 / Single UNIX Specification) while providing URTE-specific extensions for multi-scale digital twin management, cross-domain TRL Pull scheduling, ethical/planetary guardrail enforcement, and real-time intervention orchestration.

All interfaces, system calls, data structures, and behaviors are specified to maintain strict POSIX compliance for portability, while enabling the unique requirements of scale-invariant regeneration across cellular, ecological, and planetary domains.
\end{abstract}

\tableofcontents
\newpage

%==============================================================================
\chapter{Introduction and Scope}
%==============================================================================

\section{Purpose}

This specification defines the architecture, interfaces, data structures, and behavioral requirements of the \textbf{URTE Kernel} — the core operating system component that provides:

- Process and thread management for regeneration agents and digital twin simulators
- Protected address spaces for multi-scale state vectors ($u_s$)
- Inter-process communication (IPC) mechanisms optimized for cross-scale data flow
- Real-time scheduling with TRL Pull priority inheritance
- Security and guardrail enforcement as kernel policy
- POSIX-compliant file, memory, and device abstractions for clinical, ecological, and space deployment environments

\section{POSIX Compliance Statement}

The URTE Kernel \textbf{shall} be compliant with:

- \textbf{IEEE Std 1003.1-2024} (POSIX.1)
- \textbf{Single UNIX Specification, Version 4 (SUSv4)}
- \textbf{POSIX Realtime Extensions} (1003.1b)
- \textbf{POSIX Threads} (1003.1c)
- \textbf{POSIX Message Queues, Semaphores, Shared Memory} (1003.1b)

\textbf{Compliance Strategy:}
\begin{itemize}
    \item Core kernel implements a minimal POSIX-compatible interface (system calls, signals, file descriptors, processes/threads).
    \item URTE-specific functionality is provided via:
          \begin{itemize}
              \item Standard POSIX mechanisms (shared memory, message queues, real-time signals)
              \item Well-defined user-space libraries (liburte) that do not modify kernel behavior for POSIX processes
              \item Optional kernel modules loaded only when URTE-specific features are enabled (never breaking POSIX semantics for standard processes)
          \end{itemize}
\end{itemize}

The kernel passes the POSIX conformance test suite (e.g., VSX-PCTS or equivalent) when operating in standard mode.

\section{Document Scope}

\begin{itemize}
    \item Kernel architecture and layering
    \item Process model and scheduling (including TRL Pull extensions)
    \item Memory management for multi-scale digital twins
    \item IPC mechanisms for cross-layer and cross-scale communication
    \item Security model and guardrail enforcement
    \item File system abstractions (for models, telemetry, audit logs)
    \item System call interface (POSIX + URTE extensions)
    \item Boot, initialization, and shutdown
    \item ASCII architecture schemas and data structure definitions
\end{itemize}

%==============================================================================
\chapter{Kernel Architecture Overview}
%==============================================================================

\section{High-Level Kernel Layering (ASCII Schema)}

\begin{lstlisting}[style=kernelschema, caption={URTE Kernel High-Level Architecture}]
+-----------------------------------------------------------------------+
|                        USER SPACE (POSIX Processes)                   |
|  +------------------+  +------------------+  +------------------+     |
|  | Regen Agent      |  | Digital Twin     |  | Ethics Board     |     |
|  | (Clinical/Eco)   |  | Simulator        |  | Daemon           |     |
|  +------------------+  +------------------+  +------------------+     |
|  | liburte (URTE extensions) |  POSIX C Library (glibc/uclibc) |     |
+-----------------------------------------------------------------------+
                                  | System Calls
+-----------------------------------------------------------------------+
|                           URTE KERNEL                                 |
|-----------------------------------------------------------------------|
|  Kernel Mode                                                          |
|                                                                       |
|  +------------------+  +------------------+  +------------------+     |
|  | Process Manager  |  | Memory Manager   |  | Scheduler (RT +    |     |
|  | (POSIX + TRL     |  | (Multi-scale     |  | TRL Pull)          |     |
|  | Pull extensions) |  | State Vectors)   |  |                    |     |
|  +------------------+  +------------------+  +------------------+     |
|                                                                       |
|  +------------------+  +------------------+  +------------------+     |
|  | IPC Subsystem    |  | Security/        |  | VFS + Device       |     |
|  | (POSIX MQ/SHM/   |  | Guardrail        |  | Drivers            |     |
|  | Semaphores +     |  | Enforcement      |  | (Bionanorobot,     |     |
|  | URTE Channels)   |  | Module           |  | IoT, Telemetry)    |     |
|  +------------------+  +------------------+  +------------------+     |
|                                                                       |
|  +------------------+  +------------------+                           |
|  | Timer / Clock    |  | Audit & Logging  |                           |
|  | (POSIX realtime) |  | Subsystem        |                           |
|  +------------------+  +------------------+                           |
+-----------------------------------------------------------------------+
                                  | Hardware Abstraction Layer (HAL)
+-----------------------------------------------------------------------+
|                     HARDWARE / VIRTUALIZATION LAYER                   |
|  Clinical Servers | Edge Devices | IoT Swarms | Satellites | ISRU HW   |
+-----------------------------------------------------------------------+
\end{lstlisting}

\section{Core Design Principles}

\begin{enumerate}
    \item \textbf{POSIX First}: All standard POSIX interfaces behave exactly as specified. URTE extensions are additive and optional.
    \item \textbf{Scale-Invariant Abstraction}: Kernel primitives (processes, memory regions, IPC) are designed to represent entities at any scale (molecular agent, tissue simulator, ecosystem manager, planetary coordinator).
    \item \textbf{Guardrails by Design}: Ethical, planetary-boundary, and safety constraints are enforced at the kernel level (via the Security/Guardrail Enforcement Module) and cannot be bypassed by user-space processes without explicit authorization.
    \item \textbf{TRL Pull Scheduling}: The scheduler provides priority inheritance and bandwidth reservation mechanisms that accelerate lower-TRL modules when higher-TRL (clinical) modules generate revenue/data.
    \item \textbf{Real-time + Best-Effort Hybrid}: Hard real-time for safety-critical intervention loops; best-effort POSIX scheduling for simulation and learning workloads.
\end{enumerate}

%==============================================================================
\chapter{Process and Thread Model}
%==============================================================================

\section{POSIX Process Model Extensions}

The URTE Kernel supports standard POSIX processes and threads (pthreads) with the following URTE-specific process types:

\begin{itemize}
    \item \textbf{Regen Agent Process}: Long-lived or ephemeral processes representing biological or ecological intervention entities. Each has an associated \texttt{scale\_level} attribute (Molecular | Cellular | Tissue | Ecosystem | Planetary).
    \item \textbf{Digital Twin Simulator Process}: Compute-intensive processes running PINN models. Can share memory regions with other simulators for cross-scale coupling.
    \item \textbf{Guardrail Monitor Process}: Privileged processes that subscribe to kernel events and can trigger intervention pauses.
\end{itemize}

\section{Process Control Block (PCB) Extension (ASCII Schema)}

\begin{lstlisting}[style=kernelschema, caption={Extended Process Control Block (PCB) for URTE}]
struct urte_pcb {
    /* POSIX standard fields */
    pid_t           pid;
    pid_t           ppid;
    uid_t           uid, euid, suid;
    gid_t           gid, egid, sgid;
    struct task_struct *threads;        /* POSIX threads list */
    struct mm_struct   *mm;             /* Address space */
    struct files_struct *files;         /* File descriptor table */
    struct signal_struct *signal;

    /* URTE-specific extensions (stored in kernel object) */
    uint32_t        scale_level;        /* 0=Molecular ... 6=Planetary */
    uint64_t        trl_level;          /* Current TRL of this agent */
    struct urte_guardrail_policy *policy; /* Pointer to active guardrails */
    struct urte_state_vector *state;    /* Associated multi-scale state */
    bool            is_trl_pull_source; /* Clinical revenue generator? */
    struct list_head trl_pull_dependents; /* Lower-TRL modules accelerated */
    /* ... */
};
\end{lstlisting}

%==============================================================================
\chapter{Memory Management}
%==============================================================================

\section{Multi-Scale State Vector Memory Management}

The kernel provides specialized support for \textbf{multi-scale state vectors} ($u_s$) used by digital twins.

\begin{lstlisting}[style=kernelschema, caption={Multi-Scale State Vector Memory Region}]
struct urte_state_vector {
    enum scale_level level;
    size_t           dimension;         /* Size of state vector at this scale */
    void            *data;              /* Pointer to state data (in kernel or shared mem) */
    struct urte_state_vector *parent;   /* Cross-scale coupling: coarser scale */
    struct urte_state_vector *child;    /* Finer scale */
    spinlock_t       lock;              /* For concurrent access */
    struct list_head subscribers;       /* Processes monitoring this vector */
    /* Resilience & tipping point metadata */
    double           lambda_dom;
    double           variance;
    double           lag1_autocorr;
};
\end{lstlisting}

\textbf{POSIX Compliance}: These structures are exposed to user space via POSIX shared memory (\texttt{shm\_open}, \texttt{mmap}) or memory-mapped files. Standard POSIX processes see them as ordinary shared memory regions.

%==============================================================================
\chapter{Inter-Process Communication (IPC)}
%==============================================================================

\section{POSIX IPC + URTE Channels}

The kernel implements full POSIX IPC:

- Message Queues (\texttt{mq\_open}, \texttt{mq\_send}, etc.)
- Shared Memory (\texttt{shm\_open}, \texttt{mmap})
- Semaphores (\texttt{sem\_open}, \texttt{sem\_post})
- Real-time Signals

\textbf{URTE Extension}: \texttt{URTE Channels} — typed, scale-aware message channels with built-in guardrail checking on every message.

\begin{lstlisting}[style=kernelschema, caption={URTE Channel Creation (User-space API)}]
/* POSIX-compliant wrapper + URTE extension */
int urte_channel_open(const char *name, int oflag, 
                      struct urte_channel_attr *attr);
/* attr contains: scale_level, required_trl, guardrail_mask */
\end{lstlisting}

%==============================================================================
\chapter{Scheduling and TRL Pull}
%==============================================================================

\section{Scheduler Architecture}

The URTE Kernel uses a hybrid scheduler:

- \textbf{POSIX SCHED\_FIFO / SCHED\_RR} for hard real-time intervention loops
- \textbf{POSIX SCHED\_OTHER} (CFS-like) for best-effort simulation and learning
- \textbf{URTE TRL Pull Scheduler Extension}: Dynamic priority boosting and bandwidth reservation for processes that contribute to TRL advancement of dependent modules.

\begin{lstlisting}[style=kernelschema, caption={TRL Pull Scheduling Policy (Conceptual)}]
/* When a high-TRL clinical process generates revenue/data: */
if (process->is_trl_pull_source) {
    for_each_dependent(lower_trl_process) {
        boost_priority(lower_trl_process, 
                       calculate_trl_pull_boost(revenue_generated));
        reserve_bandwidth(lower_trl_process, 
                          trl_pull_bandwidth_allocation);
    }
}
\end{lstlisting}

%==============================================================================
\chapter{Security and Guardrail Enforcement}
%==============================================================================

\section{Kernel Guardrail Module}

The \textbf{Security/Guardrail Enforcement Module} is a mandatory kernel component that intercepts critical operations and enforces:

- Planetary boundary constraints (encoded as immutable kernel policies)
- Ethical veto from authorized processes (Ethics Board)
- Containment and kill-switch enforcement for bionanorobot swarms
- Audit logging of all guardrail decisions

\textbf{POSIX Compliance}: Guardrail checks are performed in addition to standard POSIX permission checks (\texttt{uid/gid}, capabilities, SELinux/AppArmor if enabled). A process that passes POSIX checks may still be denied by guardrails.

%==============================================================================
\chapter{System Call Interface}
%==============================================================================

\section{POSIX System Calls (Supported)}

All mandatory POSIX system calls are implemented, including but not limited to:

\texttt{fork}, \texttt{execve}, \texttt{exit}, \texttt{waitpid}, \texttt{kill}, \texttt{signal}, \texttt{sigaction}, \texttt{open}, \texttt{read}, \texttt{write}, \texttt{close}, \texttt{mmap}, \texttt{munmap}, \texttt{shm\_open}, \texttt{mq\_open}, \texttt{sem\_open}, \texttt{clock\_gettime}, \texttt{nanosleep}, \texttt{sched\_setscheduler}, etc.

\section{URTE-Specific System Calls (Optional Extensions)}

These system calls are provided via an optional, dynamically loadable kernel module (`urte\_core.ko`). When the module is **not loaded**, the kernel remains 100% POSIX compliant with no behavioral changes. When loaded, the extensions become available while preserving full POSIX semantics for all standard calls.

All new syscalls follow POSIX conventions:
- Return \texttt{-1} on error and set \texttt{errno}
- Use standard POSIX error codes where applicable
- Are accessible via the \texttt{syscall(2)} interface or thin \texttt{liburte} wrappers

\subsection{urte\_state\_vector\_create}

\textbf{Synopsis}
\begin{verbatim}
#include <urte/syscalls.h>

int urte_state_vector_create(struct urte_state_vector_attr *attr,
                             int flags);
\end{verbatim}

\textbf{Description}\\
Creates a new multi-scale state vector region in kernel-managed memory. The region can be mapped into the calling process (and other authorized processes) via POSIX \texttt{mmap} or returned as a file descriptor for use with \texttt{shm\_open}-style semantics.

\textbf{Arguments}
\begin{itemize}
    \item \texttt{attr} — Pointer to attributes structure defining \texttt{scale\_level}, \texttt{dimension}, initial resilience metadata, and guardrail policy reference.
    \item \texttt{flags} — Creation flags (e.g., \texttt{URTE\_SV\_SHARED}, \texttt{URTE\_SV\_GUARDRAIL\_ENFORCED}).
\end{itemize}

\textbf{Return Value}\\
On success, returns a non-negative file descriptor referring to the state vector. On error, returns \texttt{-1} and sets \texttt{errno}.

\textbf{Errors}
\begin{itemize}
    \item \texttt{EINVAL} — Invalid \texttt{scale\_level} or dimension
    \item \texttt{ENOMEM} — Insufficient kernel memory
    \item \texttt{EPERM} — Caller lacks privilege to create vectors at requested scale
    \item \texttt{EEXIST} — Vector with same name already exists (if named)
\end{itemize}

\textbf{POSIX Interaction}\\
The returned descriptor can be used with standard POSIX calls (\texttt{mmap}, \texttt{close}, \texttt{fstat}). The memory region appears as a normal anonymous or named shared memory object to POSIX processes.

\subsection{urte\_guardrail\_check}

\textbf{Synopsis}
\begin{verbatim}
int urte_guardrail_check(int sv_fd,
                         struct urte_intervention_request *req,
                         struct urte_guardrail_result *result);
\end{verbatim}

\textbf{Description}\\
Performs a synchronous guardrail evaluation on a proposed intervention against the active planetary-boundary, ethical, and safety policies associated with the state vector. This call is typically invoked by orchestration or deployment processes before physical actuation.

\textbf{Arguments}
\begin{itemize}
    \item \texttt{sv\_fd} — File descriptor of the target state vector
    \item \texttt{req} — Description of the proposed intervention (type, magnitude, location, payload)
    \item \texttt{result} — Output structure containing decision, reason code, and required actions (e.g., ``Ethics Board veto required'')
\end{itemize}

\textbf{Return Value}\\
Returns 0 on successful evaluation. The actual authorization decision is returned in \texttt{result->decision}.

\textbf{Errors}
\begin{itemize}
    \item \texttt{EBADF} — Invalid state vector descriptor
    \item \texttt{EINVAL} — Malformed request
    \item \texttt{EACCES} — Caller not authorized to query this vector's guardrails
\end{itemize}

\textbf{POSIX Interaction}\\
This call does not modify any POSIX process state. It is an additional policy check layered on top of standard access control.

\subsection{urte\_trl\_pull\_notify}

\textbf{Synopsis}
\begin{verbatim}
int urte_trl_pull_notify(pid_t source_pid,
                         uint64_t revenue_or_data_units,
                         uint32_t flags);
\end{verbatim}

\textbf{Description}\\
Notifies the kernel scheduler that a high-maturity (typically clinical) process has generated revenue or valuable longitudinal data. The kernel uses this information to apply TRL Pull priority boosting and bandwidth reservation to dependent lower-TRL modules.

\textbf{Arguments}
\begin{itemize}
    \item \texttt{source\_pid} — PID of the contributing high-TRL process
    \item \texttt{revenue\_or\_data\_units} — Quantified contribution (used for boost calculation)
    \item \texttt{flags} — Notification flags (e.g., \texttt{URTE\_TRL\_REVENUE}, \texttt{URTE\_TRL\_DATA})
\end{itemize}

\textbf{Return Value}\\
Returns 0 on success.

\textbf{Errors}
\begin{itemize}
    \item \texttt{ESRCH} — Invalid or non-URTE process
    \item \texttt{EPERM} — Caller is not the owner or authorized delegate
\end{itemize}

\textbf{POSIX Interaction}\\
This is a pure notification; it does not alter the POSIX scheduling behavior of the calling process itself. Effects are visible only on dependent processes via the extended scheduler.

\subsection{urte\_scale\_transition}

\textbf{Synopsis}
\begin{verbatim}
int urte_scale_transition(int sv_fd,
                          enum urte_scale_level new_level,
                          struct urte_transition_attr *attr);
\end{verbatim}

\textbf{Description}\\
Requests a controlled transition of a state vector (or intervention) from one scale level to another (e.g., from Tissue to Ecosystem). The kernel validates cross-scale coupling constraints and guardrails before allowing the transition.

\textbf{Arguments}
\begin{itemize}
    \item \texttt{sv\_fd} — State vector file descriptor
    \item \texttt{new\_level} — Target scale level
    \item \texttt{attr} — Transition parameters (coupling strength, validation mode, etc.)
\end{itemize}

\textbf{Return Value}\\
Returns 0 on successful transition approval.

\textbf{Errors}
\begin{itemize}
    \item \texttt{EINVAL} — Invalid scale transition requested
    \item \texttt{EPERM} — Guardrail violation or insufficient authorization
    \item \texttt{EBUSY} — Another transition is in progress on this vector
\end{itemize}

\subsection{Additional Recommended Extensions (Future)}

\begin{itemize}
    \item \texttt{urte\_intervention\_pause} — Emergency pause of physical deployment
    \item \texttt{urte\_resilience\_query} — Read current $R = -\lambda_{\rm dom}$ and tipping indicators
    \item \texttt{urte\_channel\_open} — Create typed URTE Channel with built-in guardrail attachment
\end{itemize}

These extensions are designed to be thin, auditable, and layered on top of the robust POSIX foundation. All new functionality is isolated in the optional \texttt{urte\_core} module.

%==============================================================================
\chapter{ASCII Architecture Schemas}
%==============================================================================

\section{Kernel Boot and Initialization Flow}

\begin{lstlisting}[style=kernelschema, caption={URTE Kernel Boot Sequence}]
Power On / Hypervisor Start
    |
    v
Early Boot (assembly) --> Initialize CPU, MMU, Interrupts
    |
    v
Kernel Entry (start_kernel)
    |
    +--> setup_arch()          /* URTE-specific HAL init */
    +--> sched_init()          /* POSIX + TRL Pull scheduler */
    +--> mm_init()             /* Multi-scale memory allocator */
    +--> ipc_init()            /* POSIX MQ/SHM + URTE Channels */
    +--> security_init()       /* Guardrail module load */
    +--> urte_core_init()      /* Digital twin framework, scale registry */
    |
    v
init process (PID 1) --> mount rootfs --> start urte\_init daemon
    |
    v
System ready for POSIX processes + URTE Regen Agents
\end{lstlisting}

\section{Cross-Scale Data Flow (Simplified)}

\begin{lstlisting}[style=kernelschema, caption={Cross-Scale IPC via Kernel}]
Molecular Agent (scale=0)
    | urte_channel_send(state_update)
    v
Kernel IPC Layer (guardrail check + scale translation)
    | 
    v
Tissue Simulator (scale=2)  <-- shared memory region (POSIX shm)
    |
    v
Ecosystem Manager (scale=4)
    |
    v
Planetary Coordinator (scale=6) <-- TRL Pull notification triggered
\end{lstlisting}

%==============================================================================
\chapter{Verification and Compliance}
%==============================================================================

\section{POSIX Conformance}

The kernel shall pass:
- POSIX.1-2024 Conformance Test Suite
- Real-time Profile tests (1003.1b)
- Thread safety and cancellation tests

\section{URTE-Specific Validation}

\begin{itemize}
    \item Guardrail enforcement cannot be bypassed (formal verification or extensive fault injection)
    \item TRL Pull scheduler produces measurable acceleration of dependent modules without starving POSIX best-effort processes
    \item Multi-scale state consistency across coupled simulators
\end{itemize}

%==============================================================================
\chapter{Conclusion}
%==============================================================================

The URTE POSIX-Compliant Kernel provides a solid, standards-based foundation for the Universal Regeneration Therapy Ecosystem. By strictly adhering to IEEE Std 1003.1 while adding carefully designed extensions for multi-scale regeneration, TRL Pull acceleration, and ethical guardrails, the kernel enables portable, secure, and high-performance deployment across clinical, terrestrial, and space environments.

This specification serves as the authoritative reference for kernel implementers, Capella/Arcadia modelers, and certification authorities.

\vspace{0.8cm}
\begin{center}
\textit{--- End of URTE POSIX Kernel Technical Specification v1.0 ---}
\end{center}

%==============================================================================
\appendix
%==============================================================================

\chapter{Glossary}

\begin{tabular}{ll}
\toprule
\textbf{Term} & \textbf{Definition} \\
\midrule
POSIX & Portable Operating System Interface (IEEE Std 1003.1) \\
TRL Pull & Mechanism where high-maturity modules accelerate lower-maturity ones \\
Guardrail & Kernel-enforced constraint (ethical, planetary boundary, safety) \\
Scale Level & Abstraction level from Molecular (0) to Planetary (6) \\
URTE Channel & Typed, guardrail-checked IPC primitive \\
\bottomrule
\end{tabular}

\chapter{References}

\begin{enumerate}
    \item IEEE Std 1003.1-2024, \textit{Standard for Information Technology — Portable Operating System Interface (POSIX)}
    \item The Open Group, \textit{Single UNIX Specification, Version 4}
    \item URTE Architecture Document (June 19, 2026)
    \item URTE MBSE Document (June 2026)
    \item URTE Arcadia/Capella MBSE Extension (June 2026)
    \item Raissi et al. (2019) — Physics-Informed Neural Networks
    \item Richardson et al. (2023) — Planetary Boundaries
\end{enumerate}

\end{document}
